\section{Problem formulation}\label{sec:problem}

\subsection{The prediction problem}

A secondary cone crusher is surveyed at $n$ operating points $\mathbf x_i=(c_i,f_{80,i})$, the closed-side setting $c$ and
the feed size $f_{80}$ in mm, and each survey records the throughput in t/h and the power in kW. For each output $t$ the plant
has a mean function $\mu_t(\mathbf x)$ that no one observes, and a survey returns $y_{t,i}=\mu_t(\mathbf x_i)+\varepsilon_{t,i}$,
with noise of mean zero and variance $\sigma_t^2$, independent between surveys. The noise includes the measurement error of
the survey and the variation of the ore that the two operating variables do not capture. A predictor is a rule that maps
the $n$ surveys to a prediction $\hat y_t(\mathbf x_h)$ at a new operating point $\mathbf x_h$, which may lie inside the
range of the surveys or outside it. Its quality at $\mathbf x_h$ is the expected squared error
\begin{equation}
R(\mathbf x_h)=\mathbb E\big[(y_{t,h}-\hat y_t(\mathbf x_h))^2\big],
\label{eq:risk}
\end{equation}
over the noise of the training surveys and of the new survey, with the design held fixed; averaged over a distribution of
new operating points it is the prediction risk of the predictor on that operating range. A band predictor adds an interval
$[\hat y_t-q_h,\hat y_t+q_h]$, and its quality is the probability that the interval covers the new survey and its width.

The question of the paper is, for a given plant, operating range and number of surveys, the sign and the size of the
difference of~\eqref{eq:risk} between the calibrated predictor of Section~\ref{sec:method} and three empirical predictors:
least squares on the three regressors of the discrepancy (an intercept, the standardized setting and the standardized feed
size), least squares on the intercept and the setting alone (the setting-only line, the comparator that the ten surveys could
not separate from the method~\cite{cb002}), and the training mean. It also asks how often the calibrated predictor's 80\pct{}
band covers a new survey, and which added measurement would teach the method the most about its scalars.

\subsection{Assumptions}

The noise is homoscedastic and, where a statement needs it, Gaussian; Theorem~\ref{thm:risk} needs only its first two
moments. The design is fixed in each statement and random in the study, where the theory is applied replicate by replicate.
The regressors are standardized with the training mean and population standard deviation, as the method does, so that the
discrepancy regressors of the training surveys are always centred; at a new point they are standardized with the same
training statistics. The design matrix has full column rank, which holds whenever the surveys vary both the setting and the
feed size not in exact proportion. The plant's mean functions are smooth on the operating range; they are not assumed to be the
engine at any value of its scalars. The scalars of the engine are either given (Section~\ref{sec:theory}) or learned
(Sections~\ref{sec:method-linear} and~\ref{sec:results-sampler}).

\subsection{Notation}

Table~\ref{tab:notation} collects the symbols. Bold lower case is a vector over the training surveys, bold upper case a
matrix, and a subscript $h$ marks a value at the held-out operating point. The output index $t$ is dropped whenever a
statement concerns one output.

\begin{table}[!t]
\caption{Notation of the paper: the symbols of the model, of the theory and of the study, with their meaning and their unit
(t/h for throughput and kW for power where a symbol carries the unit of the output).}
\label{tab:notation}
\centering\footnotesize
\begin{tabular}{@{}>{$}l<{$}p{0.74\columnwidth}@{}}
\toprule
\text{Symbol} & Meaning (unit)\\
\midrule
\mathbf x=(c,f_{80}) & operating point: closed-side setting and feed size (mm)\\
\mu_t(\mathbf x) & the plant's mean function of output $t$ (t/h or kW)\\
\sigma_t & standard deviation of the survey noise (t/h or kW)\\
f_t(\boldsymbol\theta;\mathbf x) & the engine's output at scalars $\boldsymbol\theta$ (t/h or kW)\\
\boldsymbol\theta,\ \boldsymbol\theta_c,\ \boldsymbol\theta_0 & the eight scalars on the sampled scale; the linearization point; the virtual plants' truth\\
\tilde{\mathbf z},\ \mathbf Z,\ \mathbf G & discrepancy regressors $(1,z_c,z_f)$; their training matrix; $\mathbf G=\mathbf Z^{\top}\mathbf Z$\\
\mathbf d_t,\ \tau_t & discrepancy coefficients and their prior standard deviation (t/h or kW)\\
\lambda,\ \mathbf A_\lambda & shrinkage $\sigma^2/\tau^2$ (dimensionless); $(\mathbf G+\lambda\mathbf I)^{-1}$\\
\boldsymbol\beta,\ \mathbf r,\ r_h & the plant's training plane, its departure from it, the departure at $\mathbf x_h$ (t/h or kW)\\
\mathbf b,\ \mathbf e,\ e_h & the engine's training plane, its departure from it, the departure at $\mathbf x_h$ (t/h or kW)\\
\boldsymbol\delta,\ \tilde{\boldsymbol\delta} & the difference $\mathbf b-\boldsymbol\beta$ of the two planes; in the eigenbasis of $\mathbf G$\\
g_k,\ \nu_k & eigenvalues of $\mathbf G$; their limit per survey $g_k/n$\\
C & curvature gain $\frac1n\sum_ie_i(2r_i-e_i)$ (squared output unit)\\
\phi^2 & mean square of the plant's feed effect beyond the setting (squared output unit)\\
\rho,\ \kappa & alignment and relative size of the engine's and the plant's curvature (dimensionless)\\
\mathbf J_t,\ \mathbf j_h & the engine's Jacobian in the scalars on the training surveys; at $\mathbf x_h$\\
\boldsymbol\Lambda,\ \mathbf K_d & prior precision of the scalars; of the discrepancy coefficients\\
R(\mathbf x_h),\ \bar R & expected squared error at $\mathbf x_h$; its fixed-design average (squared output unit)\\
\alpha,\ \kappa_q & share of the engine's curvature in a virtual plant; size of an added curvature\\
n,\ n^\star & number of training surveys; the crossing size of Proposition~\ref{prop:line}\\
\bottomrule
\end{tabular}
\end{table}

\subsection{Why simulation with a known truth}

Every quantity that decides the comparison, $\mathbf r$, $\mathbf e$ and $\boldsymbol\delta$, is defined by the plant's mean
function, which the surveys see only through noise. On ten surveys with a power noise of about 39~kW, the plant's departure
from a plane cannot be separated from the noise: the pilot computation behind this study (Section~\ref{sec:repro}) found the residual of the
measured power about the plane of the regressors to have a standard deviation of \pilotPlantKw~kW, which contains both, against
\pilotEngKw~kW for the engine's departure from the same plane. A virtual
plant fixes $\mu_t$, so that each term of Theorem~\ref{thm:risk} is known, the risk can be computed exactly, and the
simulation can check the computation. The virtual plants of Section~\ref{sec:design} are built from the engine and from the
published capacity law of Rocha et al.~\cite{rocha}, so they are crushers of the kind the method is meant for, with
curvatures of known sign and size.
